\section{Explanations}
\begin{center}
{\bf A Note to the Reader}
\end{center}
\vspace{-0.2in}
Some parts of this paper are written as if I had not yet completed the
analysis of the largest cube in a tesseract.  This is because the paper was
written in two stages; the first at a time when I didn't think I would be able
to rule out (or in) the most difficult cases.  However, both stages were
completed in 1996.  Only the Table of Contents and Appendices were added in
2013, and the abstract was updated slightly to reflect work which was actually
completed in 1996.

The is a break in the paper between Part 1 and Part 2.  This probably occurred
because I realized that I would not be able to type the little squares and
circles surrounding the plus and minus signs.  This is why the rest of the
paper is handwritten.
\newpage

\renewcommand\thetheorem{\thesection.\arabic{theorem}}
\setcounter{theorem}{0}

\section{Theorems}
\label{appendix:theorems}
\begin{theorem}[Phil's Theorem]
For $m=2$, $n=3$, if $\pm\vec{v}_1 \pm\vec{v}_2$ (with any combination of the
+ and - signs) contains one or more zeros, then $l \leq 1$ (non-optimal)
\end{theorem}
\begin{proof}
Use the normalization where the n-cube has unit side.  $\vec{v}_1 -
\vec{v}_2$, say, fits into the cube since it's a diagonal of a face.  So, its
components have absolute values $\leq 1$.  If one or more components are zero,
then the maximum squared length of $\vec{v}_1 - \vec{v}_2$ is 2 ($= 1^2 +
1^2$).  So $|\vec{v}_1 - \vec{v}_2| \leq \sqrt{2}$, and so $|\vec{v}_1| \leq
1$.
\end{proof}
\vspace{0.2in}
\begin{theorem}[GPT - Generalization of Phil's Theorem]
For $m=3$, $n=4$, if $\pm\vec{v}_1 \pm\vec{v}_2 \pm\vec{v}_3$ has 1 or more
zeros or if $\pm\vec{v}_1 \pm\vec{v}_2$ or $\pm\vec{v}_1 \pm\vec{v}_2$ or
etc., has 2 or more zeros or if $\pm\vec{v}_1$, etc. has 3 or more zeros, the
$l \leq 1$ (non-optimal).
\end{theorem}
\begin{proof}
The proof is similar to Phil's Theorem: $|\pm\vec{v}_1 \pm\vec{v}_2
\pm\vec{v}_3| \leq \sqrt{1^2+1^2+1^2} = \sqrt{3}$ so $|\vec{v}_1| \leq 1$.
$|\pm\vec{v}_1 \pm\vec{v}_2| \leq \sqrt{1^2+1^2} = \sqrt{2}$ so
$|\vec{v}_1| \leq 1$, etc. $|\pm\vec{v}_1| \leq \sqrt{1^2} = \sqrt{1}$ so
$|\vec{v}_1| \leq 1$, etc.
\end{proof}
\vspace{0.2in}
\begin{theorem}[The Zero Theorem]
For a cube in a tesseract, suppose we add a 4th row to $V$ to convert it into
a 4x4 orthogonal matrix.  Call the 4x4 matrix $W$.  If one
or more elements in the 4th row are zero, then the cube is not optimal.
\end{theorem}
\begin{proof}
Since all the elements of $W$ have absolute values $\leq 1$, it follows that
$w_{ij}^2 \leq |w_{ij}|$ for all $(i,j)$.  Assume that $w_{4j} = 0$.  Then
$R_j = \sum_{i=1}^{3} |w_{ij}| \geq \sum_{i=1}^{3} w_{ij}^2 = 1$.  Since $R_j
\geq 1$ for at least one value of $j$, $R \geq 1$ and the cube is not optimal.
\end{proof}
\newpage

\section{Additions}
On page~\pageref{refpage3}, there is an easier way
to rule out $\alpha'$, $\beta'$, and $\gamma' = 90^o$.  Since $R_1 = s\alpha'
+ c\alpha' s\beta' + c\alpha' c\beta' s\gamma'$, then $\alpha' = 90^o$ implies
$R_1 \geq 1$.  If $s\beta' = 1$ then $R_2 = c\alpha' + s\alpha' +s\alpha'
c\beta' s\gamma' \geq c\alpha' + s\alpha' \geq 1$.  If $s\gamma' = 1$ the $R_3
= c\beta' + s\beta' \geq 1$.  This method seems to work even in higher
dimensions for an m-cube in an (m+1)-cube.

Here is yet another way to prove that the minimum $R$ does not occur on the
boundary:  On page~\pageref{refpage4}, we add a fourth row to the matrix
parameterized by 3 angles \footnote{so as to produce a 4x4 orthogonal matrix.}.
If this 4th row is $(w~x~y~z)$ then we can take $z = s\gamma$, $y = s\beta
c\gamma$, $x = s\alpha c\beta c\gamma$, and $w = -c\alpha c\beta c\gamma$. If
$s\alpha = 0$ then $x = 0$.  If $s\beta = 0$ then $y = 0$.  If $s\gamma = 0$
then $z = 0$.  If $c\alpha = 0$ then $w = 0$.  If $c\beta = 0$ then $w = x =
0$.  Finally, if $c\gamma = 0$ and $w = x = y = 0$.  In all cases, we can use
the Zero Theorem
\footnote{\label{footnote:zerotheorem}see Appendix~\ref{appendix:theorems}.},
proving that the cube is not optimal because $R \geq 1$.

On page~\pageref{refpage5}, we consider
\begin{equation*}
\begin{array}{cccc}
0 & X & X & X \\
0 & X & X & X \\
X & 0 & X & X
\end{array}
\end{equation*}
where the $X$'s may be zero or nonzero.  There is a parameterization of $V$
following that may not be correct.  However, it doesn't matter because this
case is easily ruled out by the Zero Theorem\footnote{See
footnote~\ref{footnote:zerotheorem}}, which states that if we add a fourth row
to the 3x4 matrix so as to produce a 4x4 orthogonal matrix and the fourth row
contains one or more zeros then the case is not optimal.  Here, let us call
the fourth row $(w~x~y~z)$.  Since the first two columns must be orthogonal,
$w$ or $x$ or both must be zero.

On page~\pageref{refpage6}, we decided to omit the analysis of two cases.
Here we describe them for the sake of completeness.
\newpage

{\bf Two zeros in the first column:}  Assume
that $F = K = 0$ and $R = R_1 > R_2, R_3, R_4$, then we have:
\begin{equation*}
V =
\begin{bmatrix}
A & B & C & D \\
0 & G & H & J \\
0 & L & M & P
\end{bmatrix}
\end{equation*}
\begin{equation*}
T_{mod} =
\begin{bmatrix}
aB & aC & aD & 0 & 0 & 0 \\
 G &  H &  J & 0 & 0 & 0 \\
 L &  M &  P & 0 & 0 & 0
\end{bmatrix}
\end{equation*}
with $a = \pm1$.

Suppose $a_1 a(B,C,D) + a_2(G,H,J) + a_3(L,M,P) = 0$.  Then, since $(B,C,D)$,
$(G,H,J)$, and $(L,M,P)$ are all mutually orthogonal, we can multiply by
$(B,C,D)$ to get $a_1 a(B^2+C^2+D^2) = 0$.  Since $B = C = D =
0$ would lead to a non-optimal cube, it must be that $a_1 = 0$.  Similarly,
multiplying by $(G,H,J)$ yields either $a_2 = 0$ or $G = H = J = 0$.  But the
latter case means that $\vec{v}_2 = 0$ which is impossible.  Thus, $a_2 = 0$.
Then similarly, $a_3 = 0$.  Thus $T_{mod}$ has maximum rank and the cube is
not optimal.

{\bf Two zeros in a row:}  Assume that $K = L = 0$ and $R = R_1 = R_2 > R_3,
R_4$, then we have:
\begin{equation*}
V =
\begin{bmatrix}
A & B & C & D \\
F & G & H & J \\
0 & 0 & M & P
\end{bmatrix}
\end{equation*}
\begin{equation*}
T_{mod} =
\begin{bmatrix}
 aB+fG & aC+fH & aD+fJ &   0   &   0   & 0 \\
   0   &   M   &   P   &   0   &   0   & 0 \\
-bA-gF &   0   &   0   & bC+gH & bD+gJ & 0 \\
   0   &   0   &   0   &   M   &   P   & 0
\end{bmatrix}
\end{equation*}
If $a_1 \neq 0$ and $a_2 = a_3 = a_4 = 0$ then use columns 2 and 3 of
$T_{mod}$: $a_1[a(C,D)+f(H,J)] = 0$.  If $a_1 = 0$ then we have a
contradiction.  If $a(C,D)+f(H,J) = 0$ then this is non-optimal by the GPT.

If only $a_2 \neq 0$, use columns 2 and 3 to determined either $M = P = 0$ or
$a_2 = 0$.  We know that $a_2 \neq 0$, but if $M = P = 0$ then $\vec{v}_3 = 0$
which isn't allowed.

If only $a_3 \neq 0$, use columns 4 and 5 and proceed similarly to $a_1 \neq
0$.  If only $a_4 \neq 0$, use columns 4 and 5 and process similarly to $a_2 \neq
0$.

If only $a_1$ and $a_2$ are nonzero, use columns 2 and 3: $a_1[a(C,D)+f(H,J)]
+ a_2(M,P) = 0$.  Multiply through by $(M,P)$ to get $a_2(M^2+P^2) = 0$.
Since $a_2 \neq 0$, $M = P = 0$ which isn't allowed.

If only $a_1$ and $a_3$ are nonzero, use columns 2 and 3 - it's similar to the
case where $a_1 \neq 0$.

If only $a_1$ and $a_4$ are nonzero: similar.

If only $a_2$ and $a_3$ are nonzero: similar to just $a_2$.

If only $a_2$ and $a_4$ are nonzero: similar to just $a_2$ or just $a_4$.

If only $a_3$ and $a_4$ are nonzero: use columns 4 and 5 and it's like $a_1$
and $a_2$ nonzero.

If $a_1$, $a_2$ and $a_3$ are nonzero: use columns 4 and 5 and it's like just $a_3$.

If $a_1$, $a_2$ and $a_4$ are nonzero: it's like just $a_4$ nonzero.

If $a_1$, $a_3$ and $a_4$ are nonzero: it's like just $a_1$ nonzero.

If $a_2$, $a_3$ and $a_4$ are nonzero: it's like just $a_2$ nonzero.

If $a_1$, $a_2$, $a_3$ and $a_4$ are nonzero: it's like $a_1$ and $a_2$ nonzero.

On page~\pageref{refpage7}, we are discussing the case of two zeros on a
diagonal:  $A = G = 0$ with $R = R_1 = R_2 > R_3, R_4$.  I want to show how to
do only $a_3$ and $a_4$ nonzero (i.e., $a_1 = a_2 = 0$).  Use columns 4 and 5:
$a_3(-lK,bC+lM,bD+lP) + a_4(-F,H,J) = 0 \rightarrow a_3[b(0,C,D)+l(-K,M,P)] +
a_4 (-F,H,J) = 0$.  Since
\begin{equation*}
V =
\begin{bmatrix}
0 & B & C & D \\
F & 0 & H & J \\
K & L & M & P
\end{bmatrix}
\end{equation*}
if follows that $(-F,H,J)$ is orthogonal to both $(0,C,D)$ and $(-K,M,P)$.
Thus, we can multiply through by $(-F,H,J)$ to get $a_4(F^2+H^2+J^2) = 0$.  So
again we get a contradiction.

Note:

For $a_1$, $a_2$, and $a_3$ nonzero, $a_4 = 0$, use columns 4 and 5 of $T_{mod}$
to rule it out.

For $a_1$, $a_2$, and $a_4$ nonzero, $a_3 = 0$, use columns 4 and 5 as well.

For $a_1$, $a_3$, and $a_4$ nonzero, $a_2 = 0$, use columns 2 and 3.

For $a_2$, $a_3$, and $a_4$ nonzero, $a_1 = 0$, use columns 2 and 3.
\newpage

\section{Clarifications}
\begin{center}
{\bf Snug case with no zeros}
\end{center}
\vspace{-0.1in}
The proof that the snug case with no zeros is not optimal is somewhat
fragmented, so we will clarify it here.
\begin{equation*}
V =
\begin{bmatrix}
A & B & C & D \\
F & G & H & J \\
K & L & M & P
\end{bmatrix}
\end{equation*}
On page~\pageref{refpage8}, we assume that only $J$, $M$, and $P$ are
negative.  On page~\pageref{refpage9}, we assume that some linear combination
of the rows of $T$ is zero, where $T$ is given on page~\pageref{refpage10}.
On page~\pageref{refpage11}, we find that two columns of $V$ are identical.
On page~\pageref{refpage12}, we prove the
Identical Column Theorem, which shows that $V$ is not optimal because it has
two identical columns.  On
pages~\pageref{refpage13start}-\pageref{refpage13end}, we show that for the
subcase in which two columns have the same pattern of signs so there was no
loss of generality in assuming the $J$, $M$, and $P$ were the negative
entries.

On pages~\pageref{refpage14} and \pageref{refpage15}, we consider the other
sub-case, where no two columns have the same pattern of signs, and we rule it
out, using the results from pages~\pageref{refpage16} and \pageref{refpage17}.
\newpage

\section{Corrections}
On page~\pageref{refpage18} (The Largest Square in an Odd-dimensional Cube),
I don't know why I wrote ``all equal'' above and below the matrix.  Also, when
I was working on the proof, I don't think I realized that in rotating the
maximal column with both entries nonzero with a non-maximal column with a
zero, the zero would become nonzero.  However, this doesn't affect the
validity of the proof.

On page~\pageref{refpage19}, this
\begin{itemize}
\item[]{If $w_{1j} \geq 0$ for some $j$, then $v_{ij} \geq 0$ for all $i \leq m/2$}
\item[]{If $w_{1j} \leq 0$ for some $j$, then $v_{ij} \leq 0$ for all $i \leq m/2$}
\end{itemize}
should be changed to
\begin{itemize}
\item[]{If $w_{1j} > 0$ for some $j$, then $v_{ij} \geq 0$ for all $i \leq m/2$}
\item[]{If $w_{1j} = 0$ for some $j$, then $v_{ij} =    0$ for all $i \leq m/2$}
\item[]{If $w_{1j} < 0$ for some $j$, then $v_{ij} \leq 0$ for all $i \leq m/2$}
\end{itemize}
The paragraph beginning with ``However'' is correct.  The rest of
page~\pageref{refpage20}, beginning with ``Therefore'', and continuing to
``General Infinitesimal Rotations'', should be replaced by the following:
\begin{adjustwidth}{0.35in}{}
Therefore, in the upper half of $V$, all of the matrix elements are zero
except for one in each column equal to the corresponding $w_{1j}$, and if that
$w_{1j}$ is zero, then they are all zero.  Similarly for the lower half of $V$
and $w_{2j}$.

Thus, all $v_{ij}$ can only be 0, $\pm1$, $\pm\frac{1}{4}$, and
$\pm\frac{3}{4}$.  Now the squared length must be $\frac{4n-3}{4m}$, which can
never be an integer, since the numerator is odd and the denominator is even.
But only four of the $v_{ij}$ have non-integral squared length:
$\pm\frac{3}{4}$ appears twice and $\pm\frac{1}{4}$ also appears twice.  So if
$m > 4$, some rows will contain only zeros and $\pm$ ones, with an integral
squared length.  The only hope is $m = 4$, where each row must contain one
fraction.  But then two rows will have a squared length of an integer plus
$\frac{9}{16}$, and the other two will have a squared length of an integer
plus $\frac{1}{16}$, so these cannot be equal.

This is a pity because we have not found any more values of $f(m,n)$.
However, we have proved that for $m$ even and $\geq 4$, and $n$ odd, $f(m,n)$
is {\bf strictly} less than $\sqrt{\frac{4n-3}{4m}}$.
\end{adjustwidth}

On page~\pageref{refpage24} in ``so that $v_{31}$ and $v_{33}$ must have'',
the $v_{33}$ should be $v_{32}$.

On page~\pageref{refpage25}, both references to ``columns 1 and 3'' should be
changed to ``columns 2 and 4''.

On page~\pageref{refpage26}, just before ``$(B,C,D) \bot (L,M,P)$'', it should
say ``since $\vec{v}_1 \bot \vec{v}_3$''.  Also ``$(G,H,J) \pm (L,M,P)$''
should be ``$\pm (G,H,J) \pm (L,M,P)$'' and ``$\vec{v}_2 \pm \vec{v}_3$''
should be ``$\pm \vec{v}_2 \pm \vec{v}_3$'' in all 3 places.

On page~\pageref{refpage27} in ``We have now proved that a matrix with no zeros
cannot be an optimal square in a tesseract'', ``square'' should be ``cube''.

\begin{center}
{\bf A Mystery}
\end{center}
\vspace{-0.1in}
On page~\pageref{refpage21}, it says that we omit the proof for the snug case
with exactly two zeros in the same row of $V$.  However, we do prove it
starting on page~\pageref{refpage22}.  And, it's not shown on
page~\pageref{refpage23} where that other cases are shown.

Why I did this is a mystery to me.
\newpage

\section{Other}
\begin{center}
{\bf Lower bound for $f(m,m+1)$}
\end{center}
\vspace{-0.1in}
On page~\pageref{refpage28} there is a formula for a lower bound for
$f(m,m+1)$.  Unfortunately, I no longer have my derivation of this formula and
I have not succeeded in re-deriving it.

I know that I used the idea on page~\pageref{refpage29}, showing that $f(3,4)
> 1$ and assuming that it generalizes to higher dimensions, but I just haven't
been able to get a closed-form formula like the one on
page~\pageref{refpage28}.

I also know that I used the formula $\frac{v+w}{1+\frac{vw}{c^2}}$ with $c =
1$.  This is the velocity addition formula from special relativity.  If $0
\leq v < c$ and $0 \leq w < c$, then $\frac{v+w}{1+\frac{vw}{c^2}}$ is larger
than $v$ or $w$ but less that $c$.  It's a nice substitute for the ``max''
function.
\newpage
